\section{旧锚点与前缀测度漂移}
\begin{proposition}[旧锚点引起的目标偏移]
\label{prop:stale-anchor}
固定正教师 $T$ 和几何目标 $m(q)\propto q^{1-\eta}T^\eta$，$\eta>0$。
对两个正锚点 $q_j,q_k$，
\begin{equation}
d_H(m(q_j),m(q_k))=|1-\eta|d_H(q_j,q_k).
\end{equation}
若输出 $p$ 对旧目标的拟合误差不超过 $\epsilon$，
则 $d_H(p,m(q_k))\le\epsilon+|1-\eta|d_H(q_j,q_k)$。
\end{proposition}
\begin{proof}
两个目标的对数比等于
$(1-\eta)\log(q_j/q_k)$ 加动作无关常数。
跨度消去常数，并将标量乘法转换为绝对值，得到等式。
再用 Hilbert 距离的三角不等式即得输出界。
这里的误差来自目标构造所用锚点的改变；即使精确积分，仍会存在。
采样分布与积分分布不一致导致的重要性加权问题需另行处理。
\end{proof}

\begin{proposition}[前缀测度改变的值与梯度代价]
\label{prop:prefix-drift}
设 $d,d'$ 为同一前缀空间上的概率测度，
标量或向量函数 $h$ 满足 $\sup_s\|h(s)\|\le C$。
则
\begin{equation}
\|\E_{d'}h-\E_dh\|\le2C\TV(d',d).
\end{equation}
对于条件目标 $J$，
\begin{align}
\E_{d'}J(q',T)-\E_dJ(q,T)
={}&\E_d[J(q',T)-J(q,T)]\\
&+(\E_{d'}-\E_d)J(q',T).
\end{align}
所以固定前缀测度下的进展，还需覆盖第二项才能推出总体进展。
\end{proposition}
\begin{proof}
有限前缀时，对 $\sum_s[d'(s)-d(s)]h(s)$ 用三角不等式，
得到 $C\sum_s|d'-d|=2C\TV(d',d)$；一般情形用总变差测度积分。
目标差中加减 $\E_dJ(q',T)$ 得分解。
若 $0\le J\le C$，标量项可用中心化 $J-C/2$ 将系数收紧为 $C$。
同样的论证适用于梯度，但需要独立的逐前缀梯度范数界；
目标值有界本身不提供梯度界。
\end{proof}

\section{有偏拟合与误差衰减速度}
\begin{proposition}[固定对数偏差产生倾斜稳态]
\label{prop:biased-fixed-point}
考虑几何更新
$q_{k+1,a}\propto q_{k,a}^{1-\eta}T_a^\eta e^{e_a}$，
其中 $T>0$、$e$ 固定、$0<\eta<2$。
其极限为
\begin{equation}
q_{\infty,a}=\frac{T_ae^{e_a/\eta}}{\sum_bT_be^{e_b/\eta}},
\qquad d_H(q_\infty,T)=\frac{\osc e}{\eta}.
\end{equation}
若 $1<\eta<2$ 而误差允许逐轮交替，统一输入界
$\limsup d_H(q_k,T)\le\sup_k\osc e_k/(2-\eta)$ 可以达到。
\end{proposition}
\begin{proof}
令 $y_k=C_T\log(q_k/T)$，则归一化常数被中心化消去，递推为
$y_{k+1}=(1-\eta)y_k+C_Te$。
因 $|1-\eta|<1$，
$y_k=(1-\eta)^ky_0+[1-(1-\eta)^k]C_Te/\eta$，
取极限并归一化得所述稳态和跨度。
为验证交替误差的界，取
$e_k=(-1)^{k+1}e$ 和
$y_k=(-1)^kC_Te/(2-\eta)$。
直接代入递推，两侧相等，且每步跨度为 $\osc e/(2-\eta)$。
任意初值与该特解之差按 $|1-\eta|^k$ 衰减。
固定偏差与最坏时变偏差因此有不同稳态常数。
\end{proof}

\begin{proposition}[几何衰减误差的卷积与同速因子]
\label{prop:equal-rate-errors}
设 $b_{k+1}\le c b_k+E_k$，$0<c<1$，
$0\le E_k\le B\rho^k$、$0<\rho<1$。则
\begin{equation}
b_K\le c^Kb_0+
\begin{cases}
B(c^K-\rho^K)/(c-\rho),&c\ne\rho,\\
BKc^{K-1},&c=\rho.
\end{cases}
\end{equation}
若 $c=0$，则直接有 $b_K\le E_{K-1}$，$K\ge1$。
\end{proposition}
\begin{proof}
迭代不等式得
$b_K\le c^Kb_0+B\sum_{j=0}^{K-1}c^{K-1-j}\rho^j$。
当 $c\ne\rho$，乘以 $c-\rho$ 后相邻项抵消，
得到 $c^K-\rho^K$；当 $c=\rho$，每一项均为 $c^{K-1}$，
共有 $K$ 项。该式解释了误差与精确收缩同速时出现的额外线性因子。
\end{proof}

\begin{proposition}[Mean 的小噪声局部误差底]
\label{prop:mean-noise-floor}
取几何信号、Mean 更新及 $\eta=1$，令
$b_k=d_H(q_k,T)$。若输出对精确 Mean 解的跨度误差至多为 $E$，
前文局部界写成
\begin{equation}
b_{k+1}\le h(b_k)b_k+E,\qquad
h(b)=\frac{e^bb^2}{8+e^bb^2}.
\end{equation}
若 $b_0\le2E$ 且 $h(2E)\le1/2$，则
\begin{equation}
\sup_k b_k\le2E,\qquad
\limsup_k b_k\le\frac{E}{1-h(2E)}=E+O(E^3).
\end{equation}
对趋零的局部轨迹，误差 $E_k=o(b_k^3)$ 足以保留精确更新的三次首项。
\end{proposition}
\begin{proof}
$h$ 在非负半轴非减。若 $b_k\le2E$，则
$b_{k+1}\le h(2E)2E+E\le2E$，归纳得到不变邻域。
在此邻域将 $h(b_k)$ 换为固定 $h(2E)<1$，迭代并取上极限，
得到 $E/[1-h(2E)]$。由 $h(2E)=O(E^2)$ 得最后展开。
若近似输出与精确输出的 Hilbert 距离为 $E_k$，
三角不等式的两个方向给其到教师距离之差的绝对值至多为 $E_k$，
故 $o(b_k^3)$ 不改变存在的三次首项。
以 Mean 超额损失 $\delta_k$ 提供残差证书时，
$E_k\le2\sqrt{\delta_k/q_{\min,k}}$，
所以 $\delta_k=o(q_{\min,k}b_k^6)$ 是一个充分条件。
该条件依赖原有局部轨迹与精确解证书，不声称对任意共享参数更新成立。
\end{proof}

\section{局部收缩率与全局边界}
\begin{proposition}[几何更新在稳定区间之外的两个极限]
\label{prop:geometric-unstable-limit}
固定 $q_0\ne T$、$T>0$、$\eta>2$，
令 $u_a=\log(q_{0,a}/T_a)$，最大与最小值集合分别为
$I_+$、$I_-$。精确几何迭代的偶数子列趋于 $T$ 在 $I_+$ 上的条件分布，
奇数子列趋于 $T$ 在 $I_-$ 上的条件分布。
此时 $\KL(T\Vert q_k)\to\infty$，但
$\KL(q_k\Vert T)\le\log(1/T_{\min})$ 始终有限。
\end{proposition}
\begin{proof}
闭式轨迹为 $q_k(a)\propto T_ae^{(1-\eta)^ku_a}$。
偶数时指数系数趋正无穷，除以最大指数项后，
集合外的指数趋零，集合内的比例仍为 $T_a/T_b$。
奇数时同理保留最小值集合。
$u$ 非常数，两个集合均为真子集，因此每个子列都有正教师质量被学生压到零；
$-\sum_aT_a\log q_k(a)$ 发散，而 $\sum_aT_a\log T_a$ 固定。
反向 KL 则满足
$\sum_aq_a\log q_a-\sum_aq_a\log T_a\le\log(1/T_{\min})$，
因为第一项非正。跨度发散与两个方向的 KL 有不同关系。
\end{proof}

\begin{proposition}[几何逆向 KL 没有统一严格乘法收缩]
\label{prop:no-global-kl-factor}
固定正二元教师 $T$ 与 $0<\eta<1$。
不存在仅依赖 $T,\eta$ 的 $c<1$，使所有正 $q$ 都满足
$\KL(m(q)\Vert T)\le c\KL(q\Vert T)$。
\end{proposition}
\begin{proof}
取 $q_\epsilon=(1-\epsilon,\epsilon)$。
几何目标的第二坐标与第一坐标之比为
$(\epsilon/(1-\epsilon))^{1-\eta}(T_2/T_1)^\eta\to0$，
所以 $q_\epsilon$ 与 $m(q_\epsilon)$ 都趋于 $(1,0)$。
对固定正教师，逆向 KL 在闭单纯形上连续，约定 $0\log0=0$。
两者 KL 均趋于 $\log(1/T_1)>0$，比值趋一，
与任意统一 $c<1$ 矛盾。
该反例不影响教师附近的局部比例 $(1-\eta)^2$，
也不影响全局 Hilbert 收缩。
\end{proof}
